%&pdflatex
%%%%%%%%%%%%%%%%%%%%%%%%%%%%%%%%%%%%%%%%%%%%%%%%%%%%%%%%%%%%%%%%%%%%%%
%                                                                    %
% Weak behaviour of Fourier-Jacobi series,                           %
% by                                                                 %
% J. J. Guadalupe, M. P\'erez and J. L. Varona                       %
%                                                                    %
% PUBLISHED IN:                                                      %
% J. Approx. Theory 61 (1990), 222-238                               %
%                                                                    %
% Rewrited with LaTeX2e: 1-September-2004                            %
%                                                                    %
%%%%%%%%%%%%%%%%%%%%%%%%%%%%%%%%%%%%%%%%%%%%%%%%%%%%%%%%%%%%%%%%%%%%%%

\documentclass[10pt,a4paper]{article}

\nofiles % The file .aux is not necessary

\usepackage{amsmath,amsfonts,amsthm}

\newcommand{\R}{\mathbb{R}}

\newcommand{\supp}{\operatorname{supp}}
\newcommand{\sgn}{\operatorname{sgn}}
\newcommand{\longto}{\longrightarrow}

\def\hugelimto{\mathop{\,\relbar\joinrel\relbar\joinrel%
  \longrightarrow\,}\limits}


\makeatletter
\renewcommand\section{\@startsection {section}{1}{\z@}%
                                   {-3.5ex \@plus -1ex \@minus -.2ex}%
                                   {2.3ex \@plus.2ex}%
                                   {\centering \normalfont\large\scshape}}
\makeatother

\theoremstyle{plain}
\newtheorem*{theoremMNT}{Theorem}
\newtheorem*{lemmaOne}{Lema 1}
\newtheorem*{lemmaTwo}{Lema 2}
\newtheorem*{theoOne}{Theorem 1}
\newtheorem*{theoTwo}{Theorem 2}
\newtheorem*{theoThree}{Theorem 3}
\newtheorem*{corOne}{Corollary 1}

\theoremstyle{remark}
\newtheorem*{remOne}{Remark 1}
\newtheorem*{remTwo}{Remark 2}
\newtheorem*{remThree}{Remark 3}
\newtheorem*{remReferee}{Remarks}

\title{\sffamily Weak behaviour of Fourier-Jacobi series}

\author{\scshape Jos\'{e} J. Guadalupe,\footnote{This author has been
supported in part by CAICYT Grant PB85-0338.}
\ Mario P\'{e}rez, and Juan L. Varona\\[5pt]
\textit{Departamento de Matem\'{a}ticas,}
\textit{Universidad de Zaragoza,}\\
\textit{50009 Zaragoza, Spain}
}

%% 1980 Mathematics Subject Classification:
%% \subjclass{42C15, 42A50, 46E30}

%% Key words and phrases:
% \keywords{Orthogonal expansions, Jacobi polynomials,
% weak convergence, $A_p$ weights, Lorentz spaces.}

\date{}

%%%%%%%%%%%%%%%%%%%%%%%%%%%%%%%%%%%%%%%%%%%%%%%%%%%%%%
\begin{document}
%%%%%%%%%%%%%%%%%%%%%%%%%%%%%%%%%%%%%%%%%%%%%%%%%%%%%%

%%%%%%%%%%%%%%%%%%%%%%%%%%%%%%%%%%%%%%%%%%%%%%%%%%%%%%
\maketitle
\renewcommand\thefootnote{}%
\renewcommand\footnotemark{}%
\footnote{THIS PAPER HAS BEEN PUBLISHED IN:
\textit{J.\ Approx.\ Theory}
\textbf{61}, no.~2 (1990), 222--238.}%
\footnote{Copyright \copyright\ 1990 by Academic Press, Inc.}%
%%%%%%%%%%%%%%%%%%%%%%%%%%%%%%%%%%%%%%%%%%%%%%%%%%%%%%

\begin{center}
\vspace*{-24pt}
\textit{Communicated by Paul Nevai}\\[4pt]
\textrm{Received April 11, 1988; revised July 17, 1988}\\[12pt]
DEDICATED TO THE MEMORY OF J. L. RUBIO DE FRANCIA
\end{center}

\begin{abstract}
Necessary conditions for the weak convergence of
Fourier series in orthogonal polynomials are given and it is
also shown that the partial sum operator associated with the
Jacobi series is restricted weak type, but not weak type, for
the endpoints of the mean convergence interval.
\end{abstract}

%%%%%%%%%%%%%%%%%%%%%%%%%%%%%%%%%%%%%%%%%%%%%%%%%%%%%%
\section*{Introduction}
%%%%%%%%%%%%%%%%%%%%%%%%%%%%%%%%%%%%%%%%%%%%%%%%%%%%%%

Let $d\alpha$ be a finite positive Borel measure on an interval
$I \subset \R$ such that $\supp(d\alpha)$ is an infinite set
and let $p_{n}(d\alpha)$ denote the corresponding orthonormal
polynomials. For $f \in L^{1}(d\alpha)$, $S_{n}f$ stands for
the $n$th partial sum of the orthogonal Fourier expansion of $f$
in $\{p_{n}(d\alpha)\}_{n=0}^{\infty}$, that is,
$$
S_{n}(f,x) = \sum_{k=0}^{n} a_{k}p_{k}(x),
\qquad a_{k} = a_{k}(f) = \int_{I} fp_{k}\,d\alpha.
$$

The study of the convergence of $S_{n}f$ in $L^{p}(d\alpha)$
($p \neq 2$) has been discussed for several classes of orthogonal
polynomials (c.f.\ Askey-Wainger [1], Badkov [2--4], Muckenhoupt
[11--13], Newman-Rudin [16], Pollard [17--19], Wing [24]). For
instance, in the case of Jacobi polynomials
$\{P_{n}^{(\alpha,\beta)}(x)\}_{n=0}^{\infty}$
which are orthogonal in $[-1,1]$ with respect to the weight
$w(x) = (1-x)^{\alpha}(1+x)^{\beta}$, $\alpha, \beta \geq -1/2$,
Pollard proved that $|1/p-1/2| < \min\{1/(4\alpha+4),1/(4\beta+4)\}$
is a sufficient condition for the uniform boundedness
$\|S_{n}f\|_{p,w} \leq C\|f\|_{p,w}$,
which is equivalent to the convergence in $L^{p}(w)$, $1 < p < \infty$.
Newman and Rudin showed that the previous
condition is also necessary and later Muckenhoupt extended these
results to $\alpha, \beta > -1$.

The aim of this paper is to examine the weak behaviour of the
Fourier-Jacobi expansion, that is, to study if there exists a
constant $C$, independent of $n$, $y$ and $f$, such that
$$
\int_{|S_{n}(f,x)| > y} w(x)\,dx
\leq Cy^{-p} \int_{-1}^{1} |f(x)|^{p}w(x)\,dx,
\qquad y > 0,
$$
i.e., if $S_{n}$ is uniformly bounded from $L^{p}(w)$ into $L^{p}_{*}(w)$,
$1 < p < \infty$.

By using interpolation [22, Th.~3.15, p.~197], the range of $p$'s
for which there exists convergence of $S_{n}f$ in $L^{p}(w)$ is always
an interval, named mean convergence interval. Moreover, the previous
weak inequality only can be true, beside the mean convergence
interval, in its endpoints. Since for $-1 < \alpha, \beta \leq -1/2$,
the conditions $|1/p-1/2| < \min\{1/(4\alpha+4),1/(4\beta+4)\}$
are trivial for $p \in (1,\infty)$, we suppose, by symmetry,
$\alpha \geq \beta$ and $\alpha > -1/2$. Then
the mean convergence interval is
$4(\alpha+1)/(2\alpha+3) < p < 4(\alpha+1)/(2\alpha+1)$.
For the Fourier-Legendre
expansion ($\alpha = \beta = 0$) and $p = 4$ Chanillo [5]
proved that the partial sum operator is not weak type $(4,4)$,
but is restricted weak type $(4,4)$ (and $(4/3,4/3)$, by duality),
i.e., it is weakly bounded on characteristic functions.

On the other hand, M\'{a}t\'{e}, Nevai and Totik [10] obtained, in
a general way, necessary conditions for the mean convergence
of Fourier expansions.

%% Theorem (Mate-Nevai-Totik)
\begin{theoremMNT}[M\'at\'e-Nevai-Totik]
Let $d\alpha$ be such that
$\supp(d\alpha) = [-1,1]$, $\alpha' > 0$ almost everywhere,
$U$ and $V$ nonnegative Borel measurable functions such that neither
of them vanishes almost everywhere in $[-1,1]$ and $V$ is finite
on a set with positive Lebesgue measure. If $S_{n}$ is uniformly
bounded from $L^{p}(V(x)^{p}\,d\alpha)$ into $L^{p}(U(x)^{p}\,d\alpha)$,
then
\begin{itemize}
\item[(i)] $U(x)^{p} \in L^{1}(d\alpha)$, $V(x)^{-q} \in L^{1}(d\alpha)$,
$q = p/(p-1)$,
\item[(ii)] $\displaystyle \int_{-1}^{1} U(x)^{p}\alpha'(x)^{1-p/2} (1-x^{2})^{-p/4}\,dx
< \infty$,
\item[(iii)] $\displaystyle \int_{-1}^{1} V(x)^{-q}\alpha'(x)^{1-q/2} (1-x^{2})^{-q/4}\,dx
< \infty.$
\end{itemize}
\end{theoremMNT}

When $d\alpha$ and $d\beta = U(x)^{p}\,d\alpha = V(x)^{p}\,d\alpha$
are generalized Jacobi measures, these conditions turn out to
be sufficient too [2--4].

This paper is organized as follows. In Section~1 we obtain necessary
conditions for the weak convergence. These allow us to prove
that $S_{n}$ is not weak type $(p,p)$ for $p = 4(\alpha+1)/(2\alpha+3)$.
From these conditions it follows that (i), (ii) and (iii) are necessary
not only for the mean convergence but also for the weak convergence.
We end this section giving an example which shows that they are
not sufficient. In Section~2 we prove, by using similar arguments
to~[5], that $S_{n}$ is not weak type for $p = 4(\alpha+1)/(2\alpha+1)$.
Finally, in Section~3, we obtain that the partial sum operator
is restricted weak type $(p,p)$ for both endpoints of the mean
convergence interval when $\alpha, \beta \geq -1/2$
and one of them is bigger than~$-1/2$.

%%%%%%%%%%%%%%%%%%%%%%%%%%%%%%%%%%%%%%%%%%%%%%%%%%%%%%
\section*{Section 1}
%%%%%%%%%%%%%%%%%%%%%%%%%%%%%%%%%%%%%%%%%%%%%%%%%%%%%%

Assume $\supp(d\alpha) = [-1,1]$, $\alpha' > 0$ a.e.,
and let $\{p_{n}(x)\}_{n=0}^{\infty}$ be the corresponding
orthonormal polynomials. $C$ is used to denote positive constants
not necessarily the same in each occurrence and $q = p/(p-1)$.
For $f \in L^{1}(d\alpha)$ let $S_{n}f$ denote the $n$th partial
sum of the orthogonal Fourier expansion of $f$
in $\{p_{n}(x)\}_{n=0}^{\infty}$.
We want to find necessary conditions for the weak convergence
of~$S_{n}$.

%% \textsc{Lemma} 1.
\begin{lemmaOne}
Let $U$ and $V$ be weights and let $1 < p < \infty$.
If there exists a constant $C$ such that for every
$f \in L^{p}(V^{p}\,d\alpha)$ the inequality
\begin{equation}
\|S_{n}f\|_{L^{p}_{*}(U^{p}\,d\alpha)}
\leq C\|f\|_{L^{p}(V^{p}\,d\alpha)}
\tag{1.1}
\end{equation}
holds for all integers $n \geq 0$, then
\begin{equation}
\|p_{n}\|_{L^{q}(V^{-q}\,d\alpha)}
\|p_{n}\|_{L^{p}_{*}(U^{p}\,d\alpha)} \leq C.
\tag{1.2}
\end{equation}
\end{lemmaOne}

\begin{proof}
It follows from (1.1) that
$$
\|a_{n}(f)p_{n}\|_{L^{p}_{*}(U^{p}\,d\alpha)}
= \|S_{n}f-S_{n-1}f\|_{L^{p}_{*}(U^{p}\,d\alpha)}
\leq C\|f\|_{L^{p}(V^{p}\,d\alpha)}.
$$
Thus,
$$
|a_{n}(f)| \leq C \left( \|p_{n}\|_{L^{p}_{*}(U^{p}\,d\alpha)} \right)^{-1}
\|f\|_{L^{p}(V^{p}\,d\alpha)}.
$$
Therefore, every operator
\begin{align*}
p_{n}/V^{p} : L^{p}(V^{p}\,d\alpha) &\longto \R \\
f &\longmapsto \left(\int_{-1}^{1} (p_{n}/V^{p})fV^{p}\,d\alpha \right)
= a_{n}(f)
\end{align*}
is bounded and, by duality, its norm as operator coincides with
the norm as function in $L^{q}(V^{p}\,d\alpha)$. Thus
$$
\|p_{n}/V^{p}\|_{L^{q}(V^{p}\,d\alpha)}
\leq C \left(\|p_{n}\|_{L^{p}_{*}(U^{p}\,d\alpha)} \right)^{-1}
$$
and (1.2) indeed holds.
\end{proof}



In order to prove the main theorem, we will use the following
result established by M\'{a}t\'{e}, Nevai and Totik [10, Th.~2].

%% \textsc{Lemma} 2.
\begin{lemmaTwo}
Let $\supp(d\alpha) = [-1,1]$, $\alpha' > 0$
a.e.\ in $[-1,1]$ and $0 < p \leq \infty$. There exists
a constant $C$ such that if $g$ is a Lebesgue-measurable function
in $[-1,1]$, then
$$
\|\alpha'(x)^{-1/2} (1-x^{2})^{-1/4} \|_{L^{p}(|g|^{p}\,dx)}
\leq \liminf_{n \to \infty} \|p_{n}\|_{L^{p}(|g|^{p}\,dx)}.
$$
In particular, if
$$
\liminf_{n \to \infty}\|p_{n}\|_{L^{p}(|g|^{p}\,dx)} = 0
$$
then $g = 0$ a.e.
\end{lemmaTwo}

%% \textsc{Theorem} 1.
\begin{theoOne}
Let $d\alpha$, $U$ and $V$ be as in Lemma~1. If
there exists a constant $C$ such that
$$
\|S_{n}f\|_{L^{p}_{*}(U^{p}\,d\alpha)}
\leq C\|f\|_{L^{p}(V^{p}\,d\alpha)}
$$
holds for all integers $n \geq 0$ and every $f \in L^{p}(V^{p}\,d\alpha)$,
then
\begin{equation}
U^{p}, V^{-q} \in L^{1}(d\alpha),
\tag{1.3}
\end{equation}
\begin{equation}
\alpha'(x)^{-1/2} (1-x^{2})^{-1/4} \in L^{p}_{*}(U^{p}\alpha'\,dx),
\tag{1.4}
\end{equation}
\begin{equation}
\alpha'(x)^{-1/2} (1-x^{2})^{-1/4} \in L^{q}(V^{-q}\alpha'\,dx).
\tag{1.5}
\end{equation}
\end{theoOne}

\begin{proof}
Taking $n = 0$ in Lemma~1 we obtain (1.3).

In order to proof (1.4) and (1.5), we use the following result:
$$
\|f\|_{L^{p}_{*}(dm)} \leq \sup_{E}
\frac{\|f\chi_{E}\|_{r}}{\|\chi_{E}\|_{s}}
\leq \left(\frac{p}{p-r}\right)^{1/r} \|f\|_{L^{p}_{*}(dm)},
$$
where $dm$ is a Borel measure, $0 < r < p < \infty$,
$1/s = 1/r - 1/p$ and the supremum is taken over all measurable
sets $E$ such that $0 < m(E) < \infty$ [7, Lemma~V.2.8,
p.~485]. From Lemma~2 and this inequality, it follows that
$$
\|\alpha'(x)^{-1/2} (1-x^{2})^{-1/4} \|_{L^{p}_{*}(|g|^{p}\,dx)}
\leq \liminf_{n \to \infty} \|p_{n}\|_{L^{p}_{*}(|g|^{p}\,dx)}.
$$
Now, taking $\liminf_{n \to \infty}$
in (1.2), we obtain
$$
\|\alpha'(x)^{-1/2} (1-x^{2})^{-1/4} \|_{L^{q}(V^{-q}\alpha'\,dx)}
\|\alpha'(x)^{-1/2} (1-x^{2})^{-1/4} \|_{L^{p}_{*}(U^{p}\alpha'\,dx)} \leq C.
$$
As none of this norms can vanish, we get (1.4) and (1.5).
\end{proof}

An easy consequence from Theorem~1 is

%% \textsc{Corollary} 1.
\begin{corOne}
If $d\alpha$, $U$ and $V$ are as in Lemma~1 and if
$S_{n}$ is uniformly bounded from $L^{p}(V(x)^{p}\,d\alpha)$ into
$L^{p}_{*}(U(x)^{p}\,d\alpha)$
and from $L^{q}(U(x)^{-q}\,d\alpha)$ into $L^{q}_{*}(V(x)^{-q}\,d\alpha)$,
then we have
\begin{equation}
U(x)^{p} \in L^{1}(d\alpha),\qquad V(x)^{-q} \in L^{1}(d\alpha),
\tag{1.6}
\end{equation}
\begin{equation}
\int_{-1}^{1} V(x)^{-q}\alpha'(x)^{1-q/2} (1-x^{2})^{-q/4}\,dx < \infty,
\tag{1.7}
\end{equation}
\begin{equation}
\int_{-1}^{1} U(x)^{p}\alpha'(x)^{1-p/2} (1-x^{2})^{-p/4}\,dx < \infty.
\tag{1.8}
\end{equation}
\end{corOne}

%% {\underline {Remark 1}}
\begin{remOne}
Let $S_{n}$ denote the $n$th partial sum of the Fourier-Jacobi
expansion ($\alpha'(x) = (1-x)^{\alpha}(1+x)^{\beta}$,
$\alpha, \beta \geq -1/2$, $\alpha > -1/2$)
and $U(x) = V(x) = 1$. Because of (1.3) and (1.5) in Theorem~1
we obtain the following conditions
$$
|A+1| |1/p-1/2| < (A+1)/2,
\qquad (A+1)(1/p-1/2) < 1/4,
$$
where $A$ is $\alpha$ or $\beta$, the same in each statement.
As the latter inequality is not satisfied for $p = 4(\alpha+1)/(2\alpha+3)$,
it implies that $S_{n}$ is not weak type $(p,p)$ for the lower endpoint
of the interval of mean convergence. The same happens with generalized
Jacobi polynomials.
\end{remOne}

%% {\underline {Remark 2}}.
\begin{remTwo}
The conditions (1.6), (1.7) and (1.8) are the
same that (i), (ii) and (iii) in the Introduction. These are necessary
conditions for the boundedness of $S_{n}$ from $L^{p}(V(x)^{p}\,d\alpha)$
into $L^{p}(U(x)^{p}\,d\alpha)$ or equivalently from $L^{q}(U(x)^{-q}\,d\alpha)$
into $L^{q}(V(x)^{-q}\,d\alpha)$. This points out that the conditions
obtained by M\'{a}t\'{e}, Nevai and Totik are necessary not only
for the mean convergence but also for the weak convergence.
\end{remTwo}

%% {\underline {Remark 3}}.
\begin{remThree}
Let us prove that M\'{a}t\'{e}-Nevai-Totik conditions
are not sufficient for the weak convergence. Consider the Fourier-Legendre
expansion ($d\alpha = dx$), $p = 4$, and take
\begin{align*}
U(x) &= \left|\log \left( \frac{1+x}{4} \right) \right|^{-5/8}
\left|\log \left( \frac{1-x}{4} \right) \right|^{-5/8}, \\
V(x) &= \left|\log \left( \frac{1+x}{4} \right) \right|^{-3/8}
\left|\log \left( \frac{1-x}{4} \right) \right|^{-3/8}.
\end{align*}
It is immediate that $U$ and $V$ satisfy (1.6) in Corollary~1. In
order to prove the remaining conditions (1.7) and (1.8) we will
show that the weights $U$ and $V$ satisfy even stronger ones, that
is
\begin{align}
((1-x^{2})^{\delta} U(x)^{4\delta}, (1-x^{2})^{\delta} V(x)^{4\delta})
&\in A_{4}(-1,1) \quad\text{for some }\delta > 1,
\tag{1.9}\\
((1-x^{2})^{-1} U(x)^{4}, (1-x^{2})^{-1} V(x)^{4})
&\in A_{4}(-1,1),
\tag{1.10}
\end{align}
where $A_{p}(-1,1)$ stands for the Muckenhoupt $A_{p}$ classes [9, 14],
i.e.\ $(u,v) \in A_{p}(-1,1)$, $1 < p < \infty$, if
$$
\int_{I} u(x)\,dx \left( \int_{I} v(x)^{-1/(p-1)}\,dx \right)^{p-1}
\leq C|I|^{p}
$$
for every interval $I \subset [-1,1]$, and $|I|$
denotes the Lebesgue measure of the interval~$I$. The weight $w$
belongs to $A_{1}$ if $Mw(x) \leq Cw(x)$ a.e., where $M$ denotes
the Hardy-Littlewood function.

In order to prove (1.9), by using symmetry and change of variable,
it suffices to show that
$$
(u,v) = (x^{\delta} |\log x|^{-5\delta/2}, x^{\delta} |\log x|^{-3\delta/2})
\in A_{4}(0,1/2).
$$
Let
$$
w(x) = x|\log x|^{-2} = x^{-1/2} \left(x^{-1/2} |\log x|^{2/3} \right)^{-3}
= w_{1}(x)w_{2}(x)^{-3}.
$$
It is known that $w_{1} \in A_{1}$ and it is not difficult to
prove that $w_{2} \in A_{1}$. Thus, $w \in A_{4}$ by using
the factorization theorem for $A_{p}$ weights~[6]. Now, (1.9) follows
from [15, Th.~2] and it implies that the Hilbert transform $H$
is bounded from $L^{p}(v)$ into $L^{p}(u)$, which will be used later.

On the other hand, if $u(x) = x^{-1} |\log x|^{b}$
and $v(x) = x^{-1}|\log x|^{B}$, it can be shown
that $(u,v) \in A_{p}(0,1/2)$ if and only if $-b > 1$ and
$b+1 \leq B$. Then (1.10) follows obviously.

If $p_{n}(x)$ stands for the Legendre orthonormal polynomials, the
partial sum operator can be decomposed [17] as
\begin{align}
S_{n}(f,x) &= \alpha_{n}p_{n+1}(x)
\int_{-1}^{1} \frac{p_{n}(t)-p_{n+2}(t)}{x-t}\, f(t)\,dt \notag \\
&\qquad\null + \alpha_{n}(p_{n+2}(x)-p_{n}(x))
\int_{-1}^{1} \frac{p_{n+1}(t)}{x-t}\, f(t)\,dt \notag \\
&\qquad\null - \beta_{n}p_{n+1}(x) \int_{-1}^{1} p_{n+1}(t)f(t)\,dt,
\tag{1.11}
\end{align}
where $\alpha_{n}$ and $\beta_{n}$ are bounded and
\begin{equation}
(1-x^{2})^{1/4} |p_{n}(x)| \leq C,
\qquad (1-x^{2})^{-1/4} |p_{n}(x)-p_{n+2}(x)| \leq C.
\tag{1.12}
\end{equation}
We now try to estimate the three summands of the equation (1.11).
We begin by estimating the last term using H\"{o}lder's inequality,
the first part of (1.12) and~(1.10). We have
\begin{align*}
\int_{-1}^{1} |p_{n+1}(t)f(t)|\,dt
&\leq C \left( \int_{-1}^{1} (1-t^{2})^{-1/3} V(t)^{-4/3}\,dt \right)^{3/4} \\
&\qquad \times \left(\int_{-1}^{1} |f(t)V(t)|^{4}\,dt \right)^{1/4} \\
&= C_{1}\|fV\|_{4}.
\end{align*}
Therefore
\begin{multline*}
\int_{-1}^{1} \left| p_{n+1}(x)\int_{-1}^{1} p_{n+1}(t)f(t) \,dt
\right|^{4}U(x)^{4}\,dx \\
\leq C \int_{-1}^{1} (1-x^{2})^{-1}U(x)^{4}\,dx \int_{-1}^{1} |f(t)V(t)|^{4}\,dt
= C_{1}\|fV\|_{4}^{4}.
\end{multline*}
We now estimate the middle term using (1.9) and (1.12):
\begin{align*}
&\int_{-1}^{1} \left| (p_{n+2}(x)-p_{n}(x))
\int_{-1}^{1} \frac{p_{n+1}(t)f(t)}{x-t} \,dt \right|^{4}U(x)^{4}\,dx \\
&\qquad\qquad\leq C\int_{-1}^{1} |H(p_{n+1}f,x)|^{4}(1-x^{2})U(x)^{4}\,dx \\
&\qquad\qquad\leq C_{1}\int_{-1}^{1} |p_{n+1}(x)f(x)|^{4}(1-x^{2})V(x)^{4}\,dx
\leq C_{2}\|fV\|_{4}^{4}.
\end{align*}
Finally, we will prove that the first term is not weakly bounded.
The proof is by contradiction. Let us assume that there exists
a constant $C$, independent of $n$ and $f \in L^{4}(V^{4})$, such
that
$$
\int_{|p_{n+1}(x)| |H((p_{n}-p_{n-2})f,x)| > y} U(x)^{4}\,dx
\leq Cy^{-4} \|fV\|_{4}^{4}.
$$
Then, it will be enough to construct a sequence of functions
$\{f_{n}(t)\}$ such that the constant appearing in the above inequality
grows with~$n$. A slight modification of the argument used by Chanillo
proves that $C \geq (\log n)^{3/2}$.
Therefore, the partial sum operator is not weak type $(4,4)$.
\end{remThree}

%%%%%%%%%%%%%%%%%%%%%%%%%%%%%%%%%%%%%%%%%%%%%%%%%%%%%%
\section*{Section 2}
%%%%%%%%%%%%%%%%%%%%%%%%%%%%%%%%%%%%%%%%%%%%%%%%%%%%%%

Let $S_{n}$ denote the $n$th partial sum of the Fourier-Jacobi expansion
with respect to $w(x) = (1-x)^{\alpha}(1+x)^{\beta}$, being
$\alpha \geq \beta$ and $\alpha > -1/2$. Then,
the interval of mean convergence is given by
$4(\alpha+1)/(2\alpha+3)  < p < 4(\alpha+1)/(2\alpha+1)$.
Theorem~1 works to prove that $S_{n}$ is not weak type on
$L^{p}(w)$ for $p = 4(\alpha+1)/(2\alpha+3)$,
since (1.5) is not satisfied. But it is not useful to show that
$S_{n}$ is not weak type for $p = 4(\alpha+1)/(2\alpha+1)$.
It leads us to make use of other arguments.

%% \textsc{Theorem} 2.
Let $r = 4(\alpha+1)/(2\alpha+1)$. Then,
there exists no constant $C$, independent of $n$ and $f \in L^{r}(w)$, such that
\begin{equation}
\|S_{n}f\|_{L^{r}_{*}(w)} \leq C\|f\|_{L^{r}(w)}.
\tag{2.1}
\end{equation}

\begin{proof}
In what follows assume $\supp(f) \subset [0,1]$.
Let $p_{n}(x)$ denote the orthonormal polynomials with respect to
$w(x)$ and $q_{n}(x)$ the orthonormal polynomials with respect to
$w(x)(1-x^{2})$. Moreover, if (2.1) is true, the same happens if
the left integral is only taken over the set
$\{ x \in (0,1), |S_{n}(f,x)| > y \}$. Then,
we will suppose all integrals are restricted by the condition
$x \in (0,1)$.

Pollard [18] proved that in
$$
S_{n}(f,x) = \int_{-1}^{1} f(t)K_{n}(x,t)w(t)\,dt
$$
the kernel $K_{n}(x,t)$ can be decomposed in the form
\begin{equation}
K_{n}(x,t) = r_{n}T_{1}(n,x,t) + s_{n}T_{2}(n,x,t) + s_{n}T_{3}(n,x,t)
\tag{2.2}
\end{equation}
where $r_{n}$ and $s_{n}$ are bounded and
\begin{align}
T_{1}(n,x,t) &= p_{n}(x)p_{n}(t)
\tag{2.3} \\
T_{2}(n,x,t) &= (1-t^{2}) \frac{p_{n}(x)q_{n-1}(t)}{x-t}
\tag{2.4} \\
T_{3}(n,x,t) &= T_{2}(n,t,x) = (1-x^{2}) \frac{p_{n}(t)q_{n-1}(x)}{t-x}.
\tag{2.5}
\end{align}
As $\alpha \geq -1/2$, from [23, p.~169], if $x \in [0,1]$,
we have the following estimates:
\begin{align}
|p_{n}(x)| &\leq C(1-x)^{-\alpha/2-1/4},
\tag{2.6} \\
|q_{n}(x)| &\leq C(1-x)^{-\alpha/2-3/4}.
\tag{2.7}
\end{align}
Let
$$
W_{i}(f,x) = W_{i,n}(f,x) = \int_{-1}^{1} f(t)T_{i}(n,x,t)w(t)\,dt
\qquad (i = 1, 2, 3).
$$
We try to estimate the three terms
$$
\int_{|W_{i}(f,x)| > y} w(x)\,dx
\qquad (i = 1, 2, 3).
$$

($i = 1$) By using (2.6) and H\"{o}lder's inequality, we have
\begin{align*}
|W_{1}(f,x)| &= \left| p_{n}(x) \int_{0}^{1} f(t)p_{n}(t)w(t)\,dt \right| \\
&\leq C(1-x)^{-\alpha/2-1/4}
\left(\int_{0}^{1} |f(t)|^{p}w(t)\,dt \right)^{1/p} \\
&\qquad\times \left( \int_{0}^{1} (1-t)^{q(-\alpha/2-1/4)} w(t)\,dt \right)^{1/q} \\
&= C_{1}(1-x)^{-\alpha/2-1/4} \|f\|_{p,w}
\end{align*}
if $q(-\alpha/2-1/4)+\alpha > -1$.
Applying this to $p = r$ we get
$$
\int_{|W_{1}(f,x)| > y} w(x)\,dx
\leq \int_{(1-x)^{-1/4-\alpha/2} > y/(C\|f\|_{r,w})} w(x)\,dx
\leq C_{1}y^{-r} \|f\|_{r,w}^{r}
$$
what shows that $W_{1}$ is weakly bounded.

($i = 3$) Let $H$ denote the Hilbert transform. It is well
known that $H$ is bounded from $L^{p}(u)$ into $L^{p}(u)$ if and only
if the weight $u$ belongs to $A_{p}$ [9]. By using (2.6) and (2.7),
we get
\begin{multline}
\int_{0}^{1} |W_{3}(f,x)|^{p}w(x)\,dx \notag \\
\leq C \int_{0}^{1} |H(f(t)p_{n}(t)w(t),x)|^{p} (1-x)^{(1/4-\alpha/2)p+\alpha}\,dx.
\tag{2.8}
\end{multline}
Recalling that $(1-x)^{(1/4-\alpha/2)p+\alpha} \in A_{p}(0,1)$
iff $-1 < (1/4-\alpha/2)p+\alpha < p-1$, [21], and it is verified
for $p = r$, then (2.8) and (2.6) yield
\begin{multline*}
\int_{0}^{1} |W_{3}(f,x)|^{r}w(x)\,dx \\
\leq C \int_{0}^{1} |f(x)p_{n}(x)w(x)|^{r} (1-x)^{(1/4-\alpha/2)r+\alpha}\,dx
\leq C_{1} \|f\|_{r,w}^{r}
\end{multline*}
and therefore $W_{3}$ is strongly bounded.

($i = 2$) We shall prove that there is not any constant $C$,
independent of $n$ and $f$, such that
\begin{equation}
\int_{|p_{n}(x)H(f(t)q_{n-1}(t)(1-t^{2})w(t),x)| > y} w(x)\,dx
\leq Cy^{-r} \|f\|^{r}_{r,w}.
\tag{2.9}
\end{equation}
The proof is by contradition, constructing a sequence of functions
$\{f_{m,n}(t)\}$ such that the constant $C$ appearing in the previous
inequality grows with~$m$. In order to get it we try to remove
the term $(x-t)^{-1}$ in the Hilbert transform and we need sharper forms
of (2.6) and (2.7).

Because of [23, Th.~8.21.13], if $N = n+(\alpha+\beta+1)/2$
and $\gamma = -(\alpha+1/2)\pi/2$ we have
\begin{multline*}
p_{n}(\cos\theta) = (2^{\alpha^+\beta}\pi)^{-1/2}
(\sin(\theta/2))^{-\alpha-1/2} (\cos(\theta/2))^{-\beta-1/2} \\
\times [\cos(N\theta+\gamma)+(n\sin\theta)^{-1}O(1)]
\end{multline*}
where $c/n \leq \theta \leq \pi-c/n$,
$c$ being a fixed positive number.

We will restrict our attention to $\theta \leq \pi/2$
and choose $M$ large enough and such that $M-\alpha /2$
is a positive integer.
If $M\pi/n \leq \theta \leq (M+1/8)\pi/n$, we have
\begin{align*}
N\theta+\gamma
&\leq \left(n+ \frac{\alpha+\beta+1}{2}\right)
\left(M+ \frac{1}{8} \right) \frac{\pi}{n} \\
&\qquad\qquad\null- \left(\alpha+ \frac{1}{2} \right) \frac{\pi}{2}
\hugelimto_{n \to \infty} \left(M- \frac{\alpha}{2} \right) \pi
- \frac{\pi}{8}
\end{align*}
and
$$
N\theta+\gamma
\geq \left(n+ \frac{\alpha+\beta+1}{2} \right) \frac{M\pi}{n}
- \left(\alpha+ \frac{1}{2} \right) \frac{\pi}{2}
\hugelimto_{n \to \infty} \left(M- \frac{\alpha}{2} \right) \pi
- \frac{\pi}{4}.
$$
Hence, for every $\varepsilon > 0$ there exists $n_{0}$ such that
$$
\left(M- \frac{\alpha}{2} \right) \pi - \frac{\pi}{4}
- \varepsilon \leq N\theta+\gamma
\leq \left(M- \frac{\alpha}{2} \right) \pi
- \frac{\pi}{8} + \varepsilon,
\qquad n \geq n_{0}.
$$
Therefore $|\cos(N\theta+\gamma)|$
is bounded below by a positive constant for $n$ large enough, and
the same happens to $|\cos(N\theta+\gamma)| - |(n\sin\theta)^{-1}O(1)|$.
Then $p_{n}(\cos\theta) \geq C(\sin(\theta/2))^{-\alpha-1/2}$,
and taking $x = \cos\theta$ we get
\begin{equation}
|p_{n}(x)| \geq C(1-x^{2})^{-\alpha/2-1/4}
\geq C_{1}(1-x)^{-\alpha/2-1/4} \geq C_{2}n^{\alpha+1/2}
\tag{2.10}
\end{equation}
for $n \geq n_{0}$ and
$x \in [\cos \frac{(M+1/8)\pi}{n}, \cos \frac{M\pi}{n}] = I_{n}$.

We define
$$
f_{m,n}(t) = \frac{\sgn q_{n-1}(t)}{(1-t)^{\alpha/2+1/4}}
\qquad\text{if}\quad
0 \leq t \leq \cos \left(\frac{2M\pi}{m} \right),\ m < n,
$$
and $f_{m,n}(t) = 0$ elsewhere.

Now, we are going to estimate the left side in (2.9).
If $x \in I_{n}$ and $t \in \supp(f_{m,n})$ then $0 < x-t < 1-t$.
With the aid of Lemma~2, then there exists a subsequence of $n$'s
for which
\begin{multline}
|H(f_{m,n}(t)q_{n-1}(t)(1-t)^{\alpha+1} (1+t)^{\beta+1},x)| \notag \\
= \left| \int_{0}^{\cos(2M\pi/m)}
\frac{|q_{n-1}(t)|(1-t)^{\alpha+1} (1+t)^{\beta+1}}
{(1-t)^{\alpha/2+1/4}(x-t)}\,dt \right| \geq C\log m.
\tag{2.11}
\end{multline}
Thus, this and (2.10) imply
\begin{align}
&|p_{n}(x)| |H(f_{m,n}(t)q_{n-1}(t(1-t)^{\alpha+1} (1+t)^{\beta+1},x)| \notag \\
&\qquad\qquad \geq C n^{\alpha+1/2} \log m, \qquad x \in I_{n}.
\tag{2.12}
\end{align}
On the other hand, it is easy to check up
\begin{equation}
\int_{0}^{1} |f_{m,n}(t)|^{r}(1-t)^{\alpha}(1+t)^{\beta}\,dt
\leq C \int_{0}^{\cos(2M\pi/m)} (1-t)^{-1}\,dt \leq C_{1}\log m.
\tag{2.13}
\end{equation}
Finally, assume (2.9). By using (2.13) it follows that
\begin{multline}
y^{-r} \log n \geq Cy^{-r}  \int_{0}^{1} |f_{m,n}(t)|^{r}w(t)\,dt \notag \\
\geq C_{1} \int_{|p_{n}(x)H(f_{m,n}(t)q_{n-1}(t)(1-t^{2})w(t),x)| > y} w(x)\,dx.
\tag{2.14}
\end{multline}
Choose $y = Cn^{\alpha+1/2} \log m$. As $|I_{n}| = O(n^{-2})$,
using (2.12) and (2.14) we have
\begin{multline*}
Cn^{-2-2\alpha} (\log m)^{-(2\alpha+3)/(2\alpha+1)} \\
\geq C_{1} \int_{I_{n}} (1-x)^{\alpha}(1+x)^{\beta}\,dx
\geq C_{2}|I_{n}|n^{-2\alpha}
\geq C_{3}n^{-2-2\alpha},
\end{multline*}
i.e., $(\log m)^{-(2\alpha+3)/(2\alpha^+1)} \geq C$, which is absurd.
\end{proof}

%%%%%%%%%%%%%%%%%%%%%%%%%%%%%%%%%%%%%%%%%%%%%%%%%%%%%%
\section*{Section 3}
%%%%%%%%%%%%%%%%%%%%%%%%%%%%%%%%%%%%%%%%%%%%%%%%%%%%%%

The aim of this part is to examine the restricted weak behaviour
of the $n$th partial sum of the Fourier-Jacobi expansion. Assume
$w(x) = (1-x)^{\alpha}(1+x)^{\beta}$ with $\alpha \geq
\beta \geq -1/2$ and $\alpha > -1/2$, and so the
interval of mean convergence turns out to be
$4(\alpha+1)/(2\alpha+3) < p < 4(\alpha+1)/(2\alpha+1)$.
We need some notations to establish our result.
The Lebesgue measure of any set $E \subset (-1,1)$ will be denoted $|E|$,
and $\chi_{E}(x)$ will as usual denote the characteristic
function of the set~$E$. The reader is refered to [8] or [22] for
notations and results about Lorentz spaces.

%% \textsc{Theorem} 3.
\begin{theoThree}
Let $p = 4(\alpha+1)/(2\alpha+1)$ or
$p = 4(\alpha+1)/(2\alpha+3)$. Then, there exists a constant
$C$, independent of $n$ and $E \subset (-1,1)$, such that
$$
\int_{|S_{n}(\chi_{E},x)| > y} w(x)\,dx
\leq Cy^{-p} \int_{E} w(x)\,dx.
$$
\end{theoThree}

\begin{proof}
By standard duality arguments it is enough to prove
this for $p = 4(\alpha+1)/(2\alpha+1)$. Since $\beta \geq -1/2$,
we have analogous estimates to (2.6) and (2.7)
in the interval $(-1,0)$,
\begin{align}
|p_{n}(x)| &\leq C(1-x^{2})^{-1/4} w(x)^{-1/2}, \notag \\
|q_{n}(x)| &\leq C(1-x^{2})^{-3/4} w(x)^{-1/2} ,
\qquad x \in (-1,1).
\tag{3.1}
\end{align}
We proceed as in Section~2 by making the same decomposition for
the kernel in $S_{n}(\chi_{E},x)$, and we can see that $W_{1}$
and $W_{3}$ are weakly bounded. Therefore we only need to prove
that
$$
\int_{\left| \int_{-1}^{1} \chi_{E}(t)T_{2}(n,x,t)w(t)\,dt \right| > y}
w(x)\,dx \leq Cy^{-p} \int_{E} w(x)\,dx
$$
for every measurable set $E \subset (-1,1)$. By (3.1),
\begin{multline*}
\left| \int_{-1}^{1} \chi_{E}(t)T_{2}(n,x,t)w(t)\,dt \right| \\
\leq C(1-x^{2})^{-1/4} w(x)^{-1/2} |H(\chi_{E}(t)q_{n-1}(t)(1-t^{2})w(t),x)|.
\end{multline*}
Therefore, if we denote
\begin{align*}
U &= U(E,y,n) \\
&= \{x \in (-1,1) : (1-x^{2})^{-1/4} \\
&\qquad\qquad\times w(x)^{-1/2}
|H(\chi_{E}(t)q_{n-1}(t)(1-t^{2})w(t),x)| > y\}
\end{align*}
it suffices to prove
$$
y^{p} \int_{U} w(x)\,dx \leq C \int_{E} w(x)\,dx.
$$

We can write a similar proof to that of [5].
Decompose $E = E_{1} \cup E_{2}$,
$E_{1} = E \cap [0,1)$, $E_{2} = E \cap (-1,0)$, and let
$U_{1} = U(E_{1},y,n)$, $U_{2} = U(E_{2},y,n)$. Since in
$|x| \leq 3/4$ both $(1-x)$ and $(1+x)$ have lower and upper positive
bounds and since the Hilbert transform is a bounded operator
in $L^{p}(dx)$, $1 < p < \infty$ (M. Riesz's Theorem
[20]), the following inequality holds:
$$
y^{p} \int_{U \cap \{|x| \leq 3/4\}} w(x)\,dx \leq C \int_{E} w(x)\,dx.
$$
Hence, we must show that
\begin{equation}
y^{p} \int_{U_{i} \cap \{3/4 < |x| < 1\}} w(x)\,dx
\leq C \int_{E} w(x)\,dx, \qquad i = 1, 2.
\tag{3.2}
\end{equation}
We only prove (3.2) for $i = 1$ since the estimates for $i = 2$ are
made in similar way. We begin considering $x \in (-1,-3/4)$.
As $E_{1} \subset [0,1]$, the term $(x-t)^{-1}$ in the Hilbert transform
can be dropped. By using (3.1) and H\"{o}lder's inequality for
$p = 4(\alpha+1)/(2\alpha+1)$ and $q = 4(\alpha+1)/(2\alpha+3)$,
we easily get
$$
|H(\chi_{E_{1}}(t)q_{n-1}(t)(1-t^{2})w(t),x)|
\leq C \left( \int_{E_{1}} w(t)\,dt \right)^{1/p}.
$$
Then, if $\alpha = \beta$, by using (3.1) again, we have
$$
y^{p} \int_{U_{1} \cap \{-1 < x < -3/4\}} w(x)\,dx
\leq C y^{p} \int_{B} (1+x)^{\alpha}\,dx,
$$
where $B$ is the set
$$
\left\{ x \in (-1,-3/4) : C(1+x)^{-\alpha/2-1/4}
\left(\int_{E_{1}} w(t)\,dt \right)^{1/p} > y \right\}.
$$
Then
$$
y^{p} \int_{U_{1} \cap \{-1 < x < -3/4\}} w(x)\,dx)
\leq C \int_{E_{1}} w(x)\,dx
$$
is obtained straightforward.

The previous inequality is also true when $\beta < \alpha$
and it can be easily shown taking into account that the weak
norm is smaller than the strong norm.

Consequently (3.2) will be proved if we show
\begin{equation}
y^{p} \int_{U_{1} \cap \{3/4 < x < 1\}} w(x)\,dx
\leq C \int_{E} w(x)\,dx.
\tag{3.3}
\end{equation}
Let us define the sets
\begin{align*}
A &= U_{1} \cap \{3/4 < x < 1\},\\
I_{k} &= \{x \in (0,1) : 2^{-k-1} \leq 1-x < 2^{-k}\}, \\
A_{k} &= A \cap I_{k},\qquad k \geq 2.
\end{align*}
We have
$$
A = \bigcup_{k=2}^{\infty} A_{k},\qquad
A_{k} \cap A_{j} = \emptyset \ \forall k \neq j, \qquad
\bigcup_{k=2}^{\infty} I_{k} = (3/4,1).
$$
Let us also set for $k \geq 2$
\begin{align*}
E_{11}^{(k)} &= E_{1} \cap [0,1-2^{-k+1}), \\
E_{12}^{(k)} &= E_{1} \cap [1-2^{-k+1},1-2^{-k-2}), \\
E_{13}^{(k)} &= E_{1} \cap [1-2^{-k-2},1).
\end{align*}
For every $k \geq 2$, $E_{1m}^{(k)}$ ($m = 1,2,3$) are non intersecting
sets whose union is $E_{1}$. If we denote by $\chi_{m}^{(k)}$
the characteristic function of the set
$E_{1m}^{(k)}$ ($m = 1,2,3$, $k \geq 2$), then
$$
\chi_{E_{1}} = \chi_{1}^{(k)} + \chi_{2}^{(k)}
+ \chi_{3}^{(k)},\qquad k \geq 2.
$$
Again for $k \geq 2$ and $m = 1,2,3$, let
$$
A_{k}^{(m)} = \{ x \in I_{k} : |(1-x^{2})^{-1/4} w(x)^{-1/2}
H(\chi_{m}^{(k)}(t)q_{n-1}(t)(1-t^{2})w(t),x)|  > y/3 \}.
$$
Since $A_{k} \subset A_{k}^{(1)} \cup A_{k}^{(2)} \cup A_{k}^{(3)}$, then
$$
A \subset \left[ \bigcup_{k=2}^{\infty} A_{k}^{(1)} \cup A_{k}^{(3)} \right]
\cup \left[ \bigcup_{k=2}^{\infty} A_{k}^{(2)} \right].
$$
Therefore, in order to prove (3.3) it is enough to show that
\begin{equation}
y^{p} \int_{\bigcup_{k=2}^{\infty} (A_{k}^{(1)} \cup A_{k}^{(3)})} w(x)\,dx
\leq C \int_{E} w(x)\,dx.
\tag{3.4}
\end{equation}
and
\begin{equation}
\sum_{k=2}^{\infty} y^{p} \int_{A_{k}^{(2)}} w(x)\,dx
\leq C \int_{E} w(x)\,dx.
\tag{3.5}
\end{equation}

We start with $A_{k}^{(1)}$. Since $x \in I_{k}$ and
$t \in E_{11}^{(k)}$, then $1-t \geq x-t \geq 2^{-1}(1-t)$. By
using this, (3.1) and $(\alpha+1)/q - (2\alpha+3)/4 \leq 0$, we obtain
\begin{align}
&|(1-x^{2})^{-1/4} w(x)^{-1/2} H(\chi_{1}^{(k)}(t)q_{n-1}(t)(1-t^{2})w(t),x)|
\notag \\
&\qquad\leq C(1-x)^{-\alpha/2-1/4}
\int_{-1}^{1} \chi_{1}^{(k)}(t) (1-t)^{-\alpha/2-3/4} w(t)\,dt
\notag \\
&\qquad\leq C(1-x)^{-(\alpha+1)/p}
\int_{-1}^{1} \chi_{1}^{(k)}(t) (1-t)^{-(\alpha+1)/q} w(t)\,dt.
\tag{3.6}
\end{align}
By H\"{o}lder's inequality for Lorentz spaces we have
\begin{align}
&\int_{-1}^{1} \chi_{1}^{(k)}(t) (1-t)^{-(\alpha+1)/q} w(t)\,dt \notag \\
&\qquad= \|\chi_{1}^{(k)}(t) (1-t)^{-(\alpha+1)/q}\|_{(1,1),w} \notag \\
&\qquad\leq C\|\chi_{1}^{(k)}(t)\|_{(p,1),w}
\|(1-t)^{-(\alpha+1)/q}\|_{(q,\infty),w}.
\tag{3.7}
\end{align}
Also
\begin{align}
\|\chi_{1}^{(k)}(t)\|_{(p,1),w}
&\leq C\|\chi_{1}^{(k)}(t)\|^{*}_{(p,1),w} \notag \\
&= C \left( \int_{E_{11}^{(k)}} w(x)\,dx \right)^{1/p}
\leq C \left( \int_{E} w(x)\,dx \right)^{1/p}.
\tag{3.8}
\end{align}
Finally, we get
\begin{equation}
\|(1-t)^{-(\alpha+1)/q}\|_{(q,\infty),w}
= \sup_{y > 0} y^{q} \int_{|(1-t)^{-(\alpha+1)/q}| > y} w(t)\,dt < \infty.
\tag{3.9}
\end{equation}
Using (3.6), (3.7), (3.8) and (3.9) in turn we obtain for $x \in A_{k}^{(1)}$
\begin{multline}
|(1-x^{2})^{-1/4} w(x)^{-1/2} H(\chi_{1}^{(k)}(t)q_{n-1}(t)(1-t^{2})w(t),x)|
\notag \\
\leq C(1-x)^{-(\alpha+1)/p} \left( \int_{E} w(x)\,dx \right)^{1/p}.
\tag{3.10}
\end{multline}
This argument can be made analogously for $A_{k}^{(3)}$ ($k \geq 2$)
and we have the same estimates. Thus
$$
A_{k}^{(1)} \cup A_{k}^{(3)} \subset \left\{ x \in
(0,1) : C(1-x)^{-(\alpha+1)/p} \left( \int_{E} w(x)\,dx \right)^{1/p}
> y \right\},
$$
and as the same is true for
$\bigcup_{k=2}^{\infty} (A_{k}^{(1)} \cup A_{k}^{(3)})$,
then (3.4) follows.

Let us see now what happens to $A_{k}^{(2)}$, $k \geq 2$.
As $x \in I_{k}$ we can remove the term $(1+x)^{-\beta/2-1/4}$,
and so
\begin{multline*}
A_{k}^{(2)} \subset \{ x : 2^{-k-1} < 1-x < 2^{-k},\ w(x)^{-1/p} \\
\times |H(\chi_{2}^{(k)}(t)q_{n-1}(t)(1-t^{2})w(t),x)|
> C 2^{-k(\alpha/2+1/4)+\alpha k/p} y \}.
\end{multline*}
Then, by using (3.1) and M. Riesz's Theorem and the location
of $E_{12}^{(k)}$, we obtain
$$
y^{p} \int_{A_{k}^{(2)}} w(x)\,dx \leq C \int_{E_{12}^{(k)}} w(x)\,dx.
$$
As $x$ may belong to at most three intervals of the form
$\{ 1-2^{-k+1} \leq x < 1-2^{-k-2} \}$ and the
$A_{k}^{(2)}$ are non-overlapping intervals, we have
\begin{align*}
y^{p} \int_{\bigcup_{k=2}^{\infty} A_{k}^{(2)}} w(x)\,dx
&= y^{p} \sum_{k=2}^{\infty} \int_{A_{k}^{(2)}} w(x)\,dx \\
&\leq C \sum_{k=2}^{\infty} \int_{E_{12}^{(k)}} w(x)\,dx
\leq C_{1} \int_{E} w(x)\,dx,
\end{align*}
and the theorem is shown.
\end{proof}

%% \textsc{Remarks}
\begin{remReferee}
(1) As the referee has pointed out to us,
lastly L. Colzani, S. Giulini, and G. Travaglini [25] have considered weak
type boundedness of polyhedral partial sum operators on certain
compact Lie groups, proving that weak type fails in the lower
endpoint. As a consequence, results for some Jacobi-Fourier series
can be obtained from their work.

(2) We also have a proof of Theorem~3 when either $\alpha$ or
$\beta$ are less than~$-1/2$. It can be done by modifying slightly
the proof together with the use of non-uniform estimates of Jacobi
polynomials (see~[11]).
\end{remReferee}

%%%%%%%%%%%%%%%%%%%%%%%%%%%%%%%%%%%%%%%%%%%%%%%%%%%%%%
\begin{thebibliography}{99}
%%%%%%%%%%%%%%%%%%%%%%%%%%%%%%%%%%%%%%%%%%%%%%%%%%%%%%

\frenchspacing

\bibitem{1.}
\textsc{R. Askey and S. Wainger}, Mean convergence of expansions
in Laguerre and Hermite series, \textit{Amer. J. Math.} \textbf{87} (1965),
695--708.

\bibitem{2.}
\textsc{V. M. Badkov}, Convergence in the mean of Fourier series
in orthogonal polynomials, \textit{Math. Notes} \textbf{14} (1973), 651--657.

\bibitem{3.}
\textsc{V. M. Badkov}, Convergence in the mean and almost everywhere
of Fourier series in polynomials orthogonal on an interval,
\textit{Math. USSR Sbornik} \textbf{24} (1974), 223--256.

\bibitem{4.}
\textsc{V. M. Badkov}, Approximation properties of Fourier series
in orthogonal polynomials, \textit{Russian Math. Surveys} \textbf{33} (1978),
53--117.

\bibitem{5.}
\textsc{S. Chanillo}, On the weak behaviour of partial
sums of Legendre series, \textit{Trans. Amer. Math. Soc.} \textbf{268}
(1981), 367--376.

\bibitem{6.}
\textsc{R. Coifman, P. W. Jones, and J. L. Rubio de Francia},
Constructive decomposition of B.M.O. functions and
factorization of $A_p$ weights, \textit{Proc. Amer. Math. Soc.} \textbf{87}
(1983), 675--676.

\bibitem{7.}
\textsc{J. Garc\i{\i}a Cuerva and J. L. Rubio de Francia}, ``Weighted
norm inequalities and related topics'', North Holland,
1985.

\bibitem{8.}
\textsc{R. A. Hunt}, On $L(p,q)$ spaces, \textit{Enseign. Math.}
\textbf{12} (1966), 249--275.

\bibitem{9.}
\textsc{R. A. Hunt, B. Muckenhoupt, and R. L. Wheeden}, Weighted
norm inequalities for the conjugate function and the Hilbert
transform, \textit{Trans. Amer. Math. Soc.} \textbf{176} (1973), 227--251.

\bibitem{10.}
\textsc{A. M\'at\'e, P. Nevai, and V. Totik}, Necessary conditions
for weighted mean convergence of Fourier series in orthogonal
polynomials, \textit{J. Approx. Theory} \textbf{46} (1986), 314--322.

\bibitem{11.}
\textsc{B. Muckenhoupt}, Mean convergence of Jacobi series,
\textit{Proc. Amer. Math. Soc.} \textbf{23} (1969), 306--310.

\bibitem{12.}
\textsc{B. Muckenhoupt}, Mean convergence of Hermite and Laguerre
series~I, \textit{Trans. Amer. Math. Soc.} \textbf{147} (1970), 419--431.

\bibitem{13.}
\textsc{B. Muckenhoupt}, Mean convergence of Hermite and Laguerre
series~II, \textit{Trans. Amer. Math. Soc.} \textbf{147} (1970), 433--460.

\bibitem{14.}
\textsc{B. Muckenhoupt and R. L. Wheeden}, Two weight function
norm inequalities for the Hardy-Littlewood maximal function and
the Hilbert transform, \textit{Studia Math.} \textbf{55} (1976), 227--251.

\bibitem{15.}
\textsc{C. J. Neugebauer}, Inserting $A_p$-weights, \textit{Proc.
Amer. Math. Soc.} \textbf{87} (1983), 644--648.

\bibitem{16.}
\textsc{J. Newman and W. Rudin}, Mean convergence of orthogonal
series, \textit{Proc. Amer. Math. Soc.} \textbf{3} (1952), 219--222.

\bibitem{17.}
\textsc{H. Pollard}, The mean convergence of orthogonal series~I,
\textit{Trans. Amer. Math. Soc.} \textbf{62} (1947), 387--403.

\bibitem{18.}
\textsc{H. Pollard}, The mean convergence of orthogonal series~II,
\textit{Trans. Amer. Math. Soc.} \textbf{63} (1948), 355--367.

\bibitem{19.}
\textsc{H. Pollard}, The mean convergence of orthogonal series~III,
\textit{Duke Math. J.} \textbf{16} (1949), 189--191.

\bibitem{20.}
\textsc{M. Riesz}, Sur les fonctions conjugu\'{e}es, \textit{Math. Z.}
\textbf{27} (1927), 218--244.

\bibitem{21.}
\textsc{E. M. Stein}, ``Singular integrals and differentiability
properties of functions'', Princeton Univ. Press, Cambridge,
1970.

\bibitem{22.}
\textsc{E. M. Stein and G. Weiss}, ``Introduction to
Fourier Analysis on Euclidean Spaces'', Princeton Univ.
Press, New Jersey, 1975.

\bibitem{23.}
\textsc{G. Szeg\H{o}}, ``Orthogonal Polynomials'',
Amer. Math. Soc. Colloq. Publ., Vol.~23, Amer. Math. Soc., Providence,
RI, 1959.

\bibitem{24.}
\textsc{G. M. Wing}, The mean convergence of orthogonal
series, \textit{Amer. J. Math.} \textbf{72} (1950), 792-808.

\bibitem{25.}
\textsc{L. Colzani, S. Giulini, and G. Travaglini}, ``Sharp results
for the mean summability of Fourier series on compact Lie groups'',
Universit\`{a} degli Studi di Milano, Quaderno No.~12, 1988.

\end{thebibliography}

\end{document}

